\scanpage{70}
% lines 1–4, upper rectangle
On a donc un hom de suites exactes
\diagram{\begin{tikzcd}[ampersand replacement=\&,column sep=small]
0\arrow[r]\&E_1*_{O}E_2\arrow[r,"i"]\arrow[d,"{\deleted{\unsure{\mu=\mu_{E_1,E_2}}}}"']\&
E_1\otimes E_2\arrow[r,"{x_1\otimes x_2\mathbin{\unsure{\to|-}}\psi_1(x_1)\psi_2(x_2)}"]\arrow[d,"\unsure{\mu}"']\&
O\arrow[r]\arrow[d,"\mathrm{id}"]\&0\\
0\arrow[r]\&\deleted{E}L_1\otimes L_2\arrow[r,"i"]\&E_1*E_2\arrow[r,"\psi"]\&O\arrow[r]\&0
\end{tikzcd}}
% lines 5–8, splitting
Il faut définir un $\unsure{\chi|\gamma}$-scindage de l'ext. $E_1\wedge E_2$.

Il revient au $=$ \deleted{\unsure{$T$ de $E_1\wedge E_2$|$E_1*E_2$}} de se donner
\[
E_1\mathbin{*}E_2\longrightarrow L_1\otimes L_2
\]
qui sur $L_1\otimes L_2$ se réduise à la mult. par $\chi$.
% lines 9–13, pushout and mapping property
Comme
\[
E_1*E_2=E_1\otimes E_2\mathbin{\underset{\deleted{\unsure{E_1\otimes E_2}}E_1*_{O}E_2}{+}}
(\deleted{E}L_1\otimes L_2)
\]
se donner un hom lin. de $E_1*E_2$ dans un module $M$ revient au $=$ qu'un diagramme comm
% lines 14–18, centre diagram
\diagram{\begin{tikzcd}[ampersand replacement=\&,column sep=small]
\&E_1*_{O}E_2\arrow[dl,"\mathrm{inc}"']\arrow[dr,"\mu"]\&\\
E_1\otimes E_2\arrow[dr,dashed]\arrow[ddr,"\rho=\pi"']\&\&L_1\otimes L_2\arrow[dl,dashed]\arrow[ddl,"{\sigma=\chi\,\mathrm{id}_{L_1\otimes L_2}}"]\\
\&E_1\mathbin{\unsure{*|\wedge}}E_2\arrow[d,"\pi_T"]\&\\
\&M(=E_1\otimes E_2)\quad )\&
\end{tikzcd}}
% lines 19–23, last paragraph
\deleted{\unsure{Et}} on a ici $M=L_1\otimes L_2$. Considérons
\[
\rho:E_1\otimes E_2\longrightarrow L_1\otimes L_2\qquad \rho=\pi_{E_1}\otimes\pi_{E_2}
\]
i.e.
\[
\rho(x_1\otimes x_2)=
\underbrace{(\chi x_1-\psi_1(x_1)T_1)}_{\pi_{T_1}(x_1)}\otimes
\underbrace{\bigl(\chi(x_2)-\psi_2(x_2)T_2\bigr)}_{\pi_{T_2}(\unsure{x_2|x_1})}
\]
Si $\overset{\inserted{=u_1}}{x_1}\in L_1$, on trouve \deleted{$\unsure{\chi u_1}$} (comme $\psi_1(x_1)=0$)
% archive sheet number: 69
